\documentclass{article}
\usepackage[UTF8]{ctex}
\usepackage{amsmath,amssymb,amsfonts}
\usepackage{enumitem}
\usepackage{geometry}
\geometry{a4paper, margin=1in}
% === 错题本排版优化（不参与题目标识） ===
\setlength{\parindent}{0pt}
\setlength{\parskip}{0.35em plus 0.10em minus 0.08em}
\setlength{\emergencystretch}{3em}
\setlist{itemsep=0.25em, topsep=0.40em, parsep=0pt}
\allowdisplaybreaks[2]
% === 错题本排版优化结束 ===
\begin{document}
\section*{东北大学2025年高等代数真题}
\section*{66．东北大学2025年高等代数试卷（回忆版）}
一．计算题．每题15分，共90分．\\

\subsection{题目1}
1．计算 \(n\) 阶行列式

\[
\left|\begin{array}{ccccc}
2 & 2 & & & \\
-2 & 2 & 2 & & \\
& \cdot & \cdot & \cdot & \\
& & -2 & 2 & 2 \\
& & & -2 & 2
\end{array}\right|
\]

\subsection{题目2}
2．设 \(S_{1}, S_{2}\) 分别为齐次线性方程组 \(\left\{\begin{array}{l}x-z-2 w=0 ; \\ y+2 z+w=0 .\end{array}\right.\) 和 \(\left\{\begin{array}{l}3 x+6 y+z=0 ; \\ 6 x+13 y-w=0 .\end{array}\right.\) 的解空间。

\begin{itemize}
  \item[（1）] 求两个齐次线性方程组的通解。
  \item[（2）] 求 \(S_{1}+S_{2}\) 与 \(S_{1} \cap S_{2}\) 的基与维数．
\end{itemize}

\subsection{题目3}
3．设 \(A\) 为线性空间 \(V\) 上的一个线性变换，满足

\[
\begin{gathered}
\mathcal{A}\left(\varepsilon_{1}\right)=-3 \varepsilon_{1}-a \varepsilon_{2}-90 \varepsilon_{3}, \mathcal{A}\left(\varepsilon_{2}\right)=-3 \varepsilon_{1}+15 \varepsilon_{2}-90 \varepsilon_{3}, \mathcal{A}\left(\varepsilon_{3}\right)=\varepsilon_{1}-b \varepsilon_{2}+30 \varepsilon_{3} ; \\
\mathcal{A}\left(\eta_{1}\right)=6 \eta_{1}, \mathcal{A}\left(\eta_{2}\right)=12 \eta_{2}, \mathcal{A}\left(\eta_{3}\right)=c \eta_{3} .
\end{gathered}
\]

其中 \(\varepsilon_{1}, \varepsilon_{2}, \varepsilon_{3}\) 与 \(\eta_{1}, \eta_{2}, \eta_{3}\) 分别是 \(V\) 的两组基。

\begin{itemize}
  \item[（1）] 求参数 \(a, b, c\) 的值．
  \item[（2）] 求基 \(\eta_{1}, \eta_{2}, \eta_{3}\) 到 \(\varepsilon_{1}, \varepsilon_{2}, \varepsilon_{3}\) 的过渡矩阵。
\end{itemize}

\subsection{题目4}
4．设 \(V\) 是3维欧氏空间，其一组基 \(\alpha_{1}, \alpha_{2}, \alpha_{3}\) 的度量矩阵为

\[
\left(\begin{array}{lll}
5 & a & 4 \\
a & 1 & 2 \\
4 & 2 & 5
\end{array}\right) .
\]

\begin{itemize}
  \item[（1）] 求 \(a\) 的值．
  \item[（2）] 求该欧氏空间的一组标准正交基．
\end{itemize}

\subsection{题目5}
5．设 \(A\) 是复数域上的方阵，\(A\) 的全部初等因子为 \((\lambda-2)^{2},(\lambda-2)^{2},(\lambda-2)^{2}, \lambda+3, \lambda+\) \(3,(\lambda+2 \mathrm{i})^{2},(\lambda-2 \mathrm{i})^{2}\) ．

\begin{itemize}
  \item[（1）] 求 \(A\) 的特征多项式在实数域上的标准分解式。
  \item[（2）] 求 \(A\) 的所有不变因子和所有行列式因子．
\end{itemize}

\subsection{题目6}
6．已知实二次型 \(f=a x_{1}^{2}+8 x_{2}^{2}+a x_{3}^{2}-4 x_{1} x_{2}-8 x_{1} x_{3}-4 x_{2} x_{3}\) 的符号差为 2 ．

\begin{itemize}
  \item[（1）] 求 \(a\) 的值．
  \item[（2）] 用正交线性替换的方法化二次型为标准形，并写出正交线性替换。
\end{itemize}

二．证明题。\\

\subsection{题目7}
7．设 \(A\) 是实数域上的 \(n \times n\) 矩阵。证明：\(A^{2}+A^{\mathrm{T}} A=O\) 的充要条件是 \(A\) 是实反称矩阵。

\subsection{题目8}
8．设 \(A\) 是复数域上的 \(n \times n\) 矩阵。证明：

\begin{itemize}
  \item[（1）] \(A\) 与 \(A^{T}\) 有相同的特征值．
  \item[（2）] \(A\) 与 \(A^{T}\) 的特征子空间的维数相同．
  \item[（3）] 存在可逆矩阵 \(P\) ，使得 \(P A=A^{\mathrm{T}} P\) ．
\end{itemize}

\subsection{题目9}
9．设 \(A\) 是 \(n \times n\) 矩阵，线性方程组 \(A X=\beta\) 有解．证明：

\begin{itemize}
  \item[（1）] \(A X=\beta\) 线性无关解向量的个数至多为 \(n-r(A)+1\) 。
  \item[（2）] 设 \(A\) 的特征多项式中非零根的个数为 \(k\) ，则 \(k \leq r(A)\) 。（特征值重根按重数计算）
\end{itemize}

\subsection{题目10}
10．设 \(A, B\) 均为数域 \(P\) 上的 \(n\) 阶方阵，\(B\) 为可逆矩阵。满足 \(r(E-A B)+r(E+B A)=n\) 。设\(S_{1}, S_{2}\) 分别为 \((E-A B) X=0\) 与 \((E+A B) X=0\) 的解空间．证明：

\begin{itemize}
  \item[（1）] \(r(E+B A)=r(E+A B)\) ．
  \item[（2）] \(S_{1} \oplus S_{2}=P^{n}\) ．
  \item[（3）] \(A\) 是可逆矩阵．
\end{itemize}
\end{document}
